\editionPage[\pdfbookmark[0]{Corrigés des brevets blancs}{bm-bilan-A-cor-auto-parent}]{bilan-A-cor-auto}{Brevet A · Corrigé}{Automatismes expliqués}{\hyperref[bilan-A-auto]{Énoncé p.\,\pageref*{bilan-A-auto}}}
\begin{blocEdition}{bilan-A-cor-auto}
\begin{questions}
\item $\res8$ : la multiplication est prioritaire.
\item $4/6+1/6=\res{5/6}$.
\item $\res{4{,}5\times10^5}$.
\item $\res9$, la racine carrée est positive.
\item $\res{3x-12}$.
\item $\res{x=7}$.
\item $0{,}15\times80=\res{12}$.
\item $-4+1=\res{-3}$.
\item $0{,}25\times60=\res{15\text{ min}}$.
\item La troisième valeur est $\res5$.
\item $\res{3/5}$, car il y a cinq boules équiprobables.
\item $3^2=\res{9\text{ cm}^2}$.\end{questions}
\end{blocEdition}
\editionPage{bilan-A-cor-1}{Brevet A · Corrigé 1}{Un atelier de carnets}{\hyperref[bilan-A-ex-1]{Énoncé p.\,\pageref*{bilan-A-ex-1}}}
\begin{blocEdition}{bilan-A-cor-1}
\editionEtape{1 · Partager sans reste}Le nombre de carnets divise 168 et 126. $168=2^3\times3\times7$ et $126=2\times3^2\times7$. Le plus grand diviseur commun est $2\times3\times7=42$. On prépare $\res{42\text{ carnets}}$, chacun contenant $\res4$ feuilles bleues et $\res3$ blanches.\editionEtape{2 · Recettes et solde}Les recettes sont $42\times4{,}50=189$ €. Le solde, hors autres frais, est $189-78=\res{111\text{ €}}$.\editionEtape{3 · Une comparaison de prix}$5\times0{,}9=\res{4{,}50\text{ €}}$ : les prix unitaires sont égaux. La réduction s’applique au prix de 5 €.\editionRetenir{Le maximum de lots identiques correspond au plus grand diviseur commun, pas au plus petit multiple commun.}
\end{blocEdition}
\editionPage{bilan-A-cor-2}{Brevet A · Corrigé 2}{Une voile divisée}{\hyperref[bilan-A-ex-2]{Énoncé p.\,\pageref*{bilan-A-ex-2}}}
\begin{blocEdition}{bilan-A-cor-2}
\editionEtape{1 · Le triangle rectangle}Pythagore donne $BC^2=12^2+9^2=225$, donc $\res{BC=15\text{ m}}$.\editionEtape{2 · La réduction}Les alignements et le parallélisme permettent Thalès : $AM/AB=AN/AC=MN/BC=4/12=1/3$. Ainsi $\res{AN=3\text{ m}}$ et $\res{MN=5\text{ m}}$.\editionEtape{3 · Les aires}ABC a pour aire $12\times9/2=54$ m$^2$. AMN a pour aire $4\times3/2=6$ m$^2$. Le quadrilatère a donc pour aire $\res{48\text{ m}^2}$.\editionRetenir{Le rapport des aires est $1/9$, carré du rapport des longueurs $1/3$.}\label{fig-bilan-A-cor-2}\figureGeo[corrige]{brevet-A-geo}

\end{blocEdition}
\editionPage{bilan-A-cor-3}{Brevet A · Corrigé 3}{Deux locations de matériel}{\hyperref[bilan-A-ex-3]{Énoncé p.\,\pageref*{bilan-A-ex-3}}}
\begin{blocEdition}{bilan-A-cor-3}
\editionEtape{1 · Quatre journées}A coûte $18+28=\res{46\text{ €}}$ ; B coûte $\res{40\text{ €}}$. B est moins chère de 6 €.\editionEtape{2 · Le seuil}$18+7n=10n$ donne $3n=18$, donc $\res{n=6}$ ; les deux prix valent 60 €. La différence A moins B est $18-3n$. Elle est positive avant 6 : B est moins chère. Elle est négative après 6 : A est moins chère.\editionEtape{3 · Un budget}Chez A, 10 jours coûtent 88 € et 11 jours 95 € : maximum $\res{10\text{ jours}}$. Chez B, 9 jours coûtent 90 € et 10 jours 100 € : maximum $\res{9\text{ jours}}$. Les deux tarifs augmentent avec n, ce qui justifie les maximums.\label{fig-bilan-A-cor-3}\figureGeo[corrige]{brevet-A-tarifs}

\end{blocEdition}
\editionPage{bilan-A-cor-4}{Brevet A · Corrigé 4}{Des essais et un tirage}{\hyperref[bilan-A-ex-4]{Énoncé p.\,\pageref*{bilan-A-ex-4}}}
\begin{blocEdition}{bilan-A-cor-4}
\editionEtape{1 · La série}La somme pondérée vaut $0+5+16+15=36$, donc la moyenne vaut $\res{1{,}8}$. Les effectifs cumulés sont 2, 7, 15, 20 ; les rangs 10 et 11 valent 2 : médiane $\res2$.\editionEtape{2 · Deux rouges}La remise maintient la composition du sac. Le chemin RR a pour poids $(3/4)^2=\res{9/16}$.\editionEtape{3 · Exactement une bleue}Les chemins RB et BR valent chacun $3/4\times1/4=3/16$. Leur somme est $\res{6/16=3/8}$. Elle est inférieure à $9/16$, de $3/16$. Le chemin BB est exclu par « exactement une ». \label{fig-bilan-A-cor-4}\figureGeo[corrige]{brevet-A-arbre}

\end{blocEdition}
\editionPage{bilan-B-cor-auto}{Brevet B · Corrigé}{Automatismes expliqués}{\hyperref[bilan-B-auto]{Énoncé p.\,\pageref*{bilan-B-auto}}}
\begin{blocEdition}{bilan-B-cor-auto}
\begin{questions}
\item $\res6$.
\item $3/5\times10/9=\res{2/3}$.
\item $\res{10^{-2}}$.
\item $\res{2^2\times3\times5}$.
\item $\res{(x-7)(x+7)}$.
\item $\res{x=-5}$.
\item $50\times1{,}12=\res{56}$.
\item $-3x=12$, donc $\res{-4}$.
\item $\res{2400\text{ L}}$.
\item $9-(-2)=\res{11}$.
\item $1-0{,}3=\res{0{,}7}$.
\item $8\times0{,}5=\res{4\text{ cm}}$.\end{questions}
\end{blocEdition}
\editionPage{bilan-B-cor-1}{Brevet B · Corrigé 1}{Un abri et sa pente}{\hyperref[bilan-B-ex-1]{Énoncé p.\,\pageref*{bilan-B-ex-1}}}
\begin{blocEdition}{bilan-B-cor-1}
\editionEtape{1 · Longueur du versant}$BC^2=4^2+2{,}5^2=22{,}25$, donc $BC=\sqrt{22{,}25}\approx\res{4{,}72\text{ m}}$.\editionEtape{2 · Angle de pente}$\tan B=2{,}5/4=0{,}625$. La fonction inverse donne $B\approx32{,}005^\circ$, soit $\res{32^\circ}$.\editionEtape{3 · Achat entier}L’aire utile vaut $6\sqrt{22{,}25}\approx28{,}3019$ m$^2$. Avec les découpes : environ 30,5660 m$^2$. Le nombre de plaques doit être au moins $30{,}5660/2\approx15{,}283$. Il faut $\res{16\text{ plaques}}$. Quinze ne couvriraient que 30 m$^2$.\label{fig-bilan-B-cor-1}\figureGeo[corrige]{brevet-B-geo}

\end{blocEdition}
\editionPage{bilan-B-cor-2}{Brevet B · Corrigé 2}{Une bouteille et un débit}{\hyperref[bilan-B-ex-2]{Énoncé p.\,\pageref*{bilan-B-ex-2}}}
\begin{blocEdition}{bilan-B-cor-2}
\editionEtape{1 · Capacité}Le volume vaut $\pi\times4^2\times20=\res{320\pi\text{ cm}^3}$. Un cm$^3$ vaut un mL : capacité $\res{1005\text{ mL}}$ au mL.\editionEtape{2 · Remplissage}Le volume visé vaut $0{,}9\times320\pi=288\pi$ mL. La durée vaut $288\pi/50\approx\res{18{,}1\text{ s}}$.\editionEtape{3 · Réduction}Le facteur de volume est $0{,}75^3=0{,}421875$. La nouvelle capacité vaut $135\pi$ mL, soit $\res{424\text{ mL}}$. Elle représente 42,1875 \% de la première ; 75 \% est le facteur des seules longueurs.\label{fig-bilan-B-cor-2}\figureGeo[corrige]{brevet-B-cylindre}

\end{blocEdition}
\editionPage{bilan-B-cor-3}{Brevet B · Corrigé 3}{Un script de figures}{\hyperref[bilan-B-ex-3]{Énoncé p.\,\pageref*{bilan-B-ex-3}}}
\begin{blocEdition}{bilan-B-cor-3}
\editionEtape{1 · Une variable à chaque tour}Les carrés utilisent successivement $\res{10,15,20,25}$ unités : l’ajout de 5 vient après le tracé.\editionEtape{2 · Longueur dessinée}La somme des côtés vaut 70. Les périmètres totalisent $4\times70=\res{280}$ unités.\editionEtape{3 · État final}Après le quatrième ajout, c vaut $\res{30}$. Les déplacements sans tracer valent 15, 20, 25, 30, donc leur somme vaut $\res{90}$ unités. La boucle s’arrête après quatre tours ; aucun appel carré(30) n’est exécuté.\label{fig-bilan-B-cor-3}\figureGeo[corrige]{brevet-B-script}

\end{blocEdition}
\editionPage{bilan-B-cor-4}{Brevet B · Corrigé 4}{Un jeu à deux étapes}{\hyperref[bilan-B-ex-4]{Énoncé p.\,\pageref*{bilan-B-ex-4}}}
\begin{blocEdition}{bilan-B-cor-4}
\editionEtape{1 · Chemins}A-rouge a pour poids $3/4\times1/2=\res{3/8}$. B-bleue a pour poids $1/4\times1/2=\res{1/8}$.\editionEtape{2 · Événement gagnant}Les chemins sont incompatibles : $P(G)=3/8+1/8=\res{1/2}$. Une série finie peut avoir un nombre différent de 50 gains ; 50 est une valeur moyenne théorique sur 100, pas une garantie.\editionEtape{3 · Bilan net}En cas de gain, le bilan net est $6-2=4$ points ; sinon il est $-2$. Le bilan moyen vaut $4\times1/2+(-2)\times1/2=\res{1\text{ point}}$. Il ne correspond pas à un résultat certain d’une partie.\label{fig-bilan-B-cor-4}\figureGeo[corrige]{brevet-B-arbre}

\end{blocEdition}
